\documentclass[a4paper,12pt]{article}

% --- Настройка локализации и кодировки ---
\usepackage[T2A]{fontenc}
\usepackage[utf8]{inputenc}
\usepackage[russian]{babel}

% --- Геометрия страницы по ГОСТ (левое 30 мм, правое 15 мм, верх/низ 20 мм) ---
\usepackage[left=30mm, right=15mm, top=20mm, bottom=20mm]{geometry}

% --- Шрифты и интервалы ---
\usepackage{times}
\usepackage{setspace}
\onehalfspacing % Полуторный междустрочный интервал
\setlength{\parindent}{1.25cm} % Абзацный отступ
\usepackage{indentfirst} % Отступ для первой строки раздела

% --- Таблицы, списки, графика ---
\usepackage{tabularx}
\usepackage{longtable}
\usepackage{array}
\usepackage{booktabs}
\usepackage{enumitem}
\setlist{nosep, leftmargin=1.5cm} % Компактные списки

% --- Настройка оглавления и заголовков ---
\usepackage{tocloft}
\renewcommand{\cftsecdotsep}{\cftdotsep} % Точки в оглавлении для разделов
\usepackage{titlesec}
\titleformat{\section}{\normalfont\bfseries}{\thesection}{1em}{}
\titleformat{\subsection}{\normalfont\bfseries}{\thesubsection}{1em}{}
\titleformat{\subsubsection}{\normalfont\bfseries}{\thesubsubsection}{1em}{}

% --- Гиперссылки ---
\usepackage[colorlinks=false, hidelinks]{hyperref}

\begin{document}

% ==============================================================================
% ТИТУЛЬНЫЙ ЛИСТ (в соответствии с Шаблон ТЗ.pdf)
% ==============================================================================
% \begin{titlepage}
%     \fontsize{10pt}{12pt}\selectfont
%     \noindent
%     \begin{minipage}[t]{0.48\textwidth}
%         \textbf{СОГЛАСОВАНО}\\
%         Преподаватель по анализу и разработке ТЗ\\[5mm]
%         \underline{\hspace{3.5cm}} Фамилия И.О.\\[2mm]
%         «\underline{\hspace{0.8cm}}» \underline{\hspace{2.5cm}} 20\underline{\hspace{0.5cm}} г.
%     \end{minipage}
%     \hfill
%     \begin{minipage}[t]{0.48\textwidth}
%         \textbf{УТВЕРЖДАЮ}\\
%         Руководитель ОП\\[5mm]
%         \underline{\hspace{3.5cm}} Фамилия И.О.\\[2mm]
%         «\underline{\hspace{0.8cm}}» \underline{\hspace{2.5cm}} 20\underline{\hspace{0.5cm}} г.
%     \end{minipage}

%     \vspace{4.5cm}

%     \begin{center}
%         \fontsize{16pt}{18pt}\selectfont
%         \textbf{ТЕХНИЧЕСКОЕ ЗАДАНИЕ}\\[4mm]
%         \fontsize{13pt}{15pt}\selectfont
%         на разработку веб-платформы для велосипедистов и\\
%         организаторов спортивных мероприятий «43Колеса»
%     \end{center}

%     \vspace{4cm}

%     \noindent
%     \begin{minipage}[t]{0.48\textwidth}
%         \textbf{СОГЛАСОВАНО}\\
%         Студент ВятГУ\\[5mm]
%         \underline{\hspace{3.5cm}} Ковальский И.И.\\[2mm]
%         «\underline{\hspace{0.8cm}}» \underline{\hspace{2.5cm}} 20\underline{\hspace{0.5cm}} г.
%     \end{minipage}
%     \hfill
%     \begin{minipage}[t]{0.48\textwidth}
%         \textbf{СОГЛАСОВАНО}\\
%         Преподаватель по веб-разработке\\[5mm]
%         \underline{\hspace{3.5cm}} Самоделкин П.А.\\[2mm]
%         «\underline{\hspace{0.8cm}}» \underline{\hspace{2.5cm}} 20\underline{\hspace{0.5cm}} г.
%     \end{minipage}

%     \vspace{1cm}

%     \begin{flushright}
%     \begin{minipage}[t]{0.48\textwidth}
%         \textbf{СОГЛАСОВАНО}\\
%         Преподаватель\\[5mm]
%         \underline{\hspace{3.5cm}} Фамилия И.О.\\[2mm]
%         «\underline{\hspace{0.8cm}}» \underline{\hspace{2.5cm}} 20\underline{\hspace{0.5cm}} г.
%     \end{minipage}
%     \end{flushright}

%     \vfill
%     \begin{center}
%         Киров\\
%         2026
%     \end{center}
% \end{titlepage}

% ==============================================================================
% СОДЕРЖАНИЕ (страница 2)
% ==============================================================================
\newpage
\setcounter{page}{2}
\tableofcontents
\newpage

% ==============================================================================
% ВВЕДЕНИЕ
% ==============================================================================
\section*{Введение}
\addcontentsline{toc}{section}{Введение}

Настоящее техническое задание (ТЗ) распространяется на разработку веб-платформы «43Колеса», представляющей собой специализированный информационный портал и сервис координации веломероприятий для любителей и профессионалов велоспорта в Кировской области и смежных регионах.

Техническое задание определяет совокупность функциональных, технических, эксплуатационных и качественных требований к программному обеспечению (ПО), регламентирует этапы проектирования, разработки, тестирования, а также порядок проведения приемо-сдаточных испытаний. Документ служит нормативным основанием для контроля соответствия разработанного программного решения требованиям заказчика и будущих пользователей.

% ==============================================================================
% 1. ТЕРМИНЫ И ОПРЕДЕЛЕНИЯ
% ==============================================================================
\section{Термины и определения}

В настоящем документе применяются следующие термины с соответствующими определениями:
\begin{itemize}
	\item \textbf{Веломероприятие (событие)} — запланированный заезд, велогонка, бревет, тренировка или любительский сбор, обладающий фиксированной датой, регламентом, стартовыми координатами и зарегистрированными участниками.
	\item \textbf{Timeline-доска} — интерактивный элемент пользовательского интерфейса, отображающий события в виде хронологической временной шкалы с фильтрацией по дисциплинам.
	\item \textbf{GPX-трек (GPS Exchange Format)} — структурированный текстовый XML-файл, содержащий географические координаты маршрута заезда, путевые точки и перепады высот.
	\item \textbf{Организатор («Коллега»)} — авторизованный пользователь, обладающий правами создания, редактирования карточек заездов и модерации участников соревнований.
	\item \textbf{Спортсмен (пользователь)} — зарегистрированный участник платформы, имеющий персональный профиль, личную статистику и возможность регистрации на заезды.
	\item \textbf{Стартовый городок} — географическая точка сбора и старта участников, отображаемая на интерактивной карте платформы.
	\item \textbf{Бревет} — квалификационный заезд на длинную дистанцию (от 200 км) с лимитом времени и контрольными пунктами по правилам Audax Club Parisien.
\end{itemize}

% ==============================================================================
% 2. ПЕРЕЧЕНЬ СОКРАЩЕНИЙ
% ==============================================================================
\section{Перечень сокращений}

В техническом задании используются следующие сокращения:
\begin{itemize}
	\item \textbf{API} (Application Programming Interface) — программный интерфейс приложения;
	\item \textbf{БД} — база данных;
	\item \textbf{ИС} — информационная система;
	\item \textbf{ПП} — программный продукт;
	\item \textbf{ПО} — программное обеспечение;
	\item \textbf{СУБД} — система управления базами данных;
	\item \textbf{ТЗ} — техническое задание;
	\item \textbf{UI} (User Interface) — пользовательский интерфейс;
	\item \textbf{UX} (User Experience) — пользовательский опыт;
	\item \textbf{CRUD} (Create, Read, Update, Delete) — базовые операции с сущностями;
	\item \textbf{JWT} (JSON Web Token) — открытый стандарт защищенной передачи данных аутентификации;
	\item \textbf{REST} (Representational State Transfer) — архитектурный стиль взаимодействия компонентов распределенной системы.
\end{itemize}

% ==============================================================================
% 3. ОСНОВНЫЕ СВЕДЕНИЯ О РАЗРАБОТКЕ
% ==============================================================================
\section{Основные сведения о разработке}

\subsection{Наименование разработки}
Полное наименование: «Программный комплекс веб-платформы для велосипедистов и организаторов спортивных заездов «43Колеса»».\\
Краткое наименование: «ИС «43Колеса»».

\subsection{Цель и задачи}
Целью разработки является создание централизованной веб-платформы, объединяющей спортивный календарь, картографию треков, метеоинформер и профили велосипедистов в рамках единой информационной среды.

Для достижения поставленной цели решаются следующие задачи:
\begin{itemize}
	\item автоматизация публикации и структурирования сведений о региональных заездах;
	\item упрощение процедуры подачи заявок и регистрации спортсменов на старты;
	\item обеспечение доступа участников к прогнозу погоды и картографическому треку;
	\item реализация оперативного информирования через чат-бота социальной сети ВКонтакте.
\end{itemize}

\subsection{Сведения об участниках разработки}
\begin{itemize}
	\item \textbf{Заказчик:} Кафедра электронных вычислительных машин ФГБОУ ВО «ВятГУ».
	\item \textbf{Разработчик:} Студент гр. ИВТб-3301-05-00 Ковальский И.И.
	\item \textbf{Руководитель разработки:} Старший преподаватель кафедры ЭВМ Самоделкин П.А.
\end{itemize}

\subsection{Сроки разработки}
Плановый период выполнения работ: февраль 2026 г. — октябрь 2026 г.

\subsection{Назначение разработки}
\subsubsection{Функциональное назначение}
Функциональным назначением ПП является предоставление пользователям веб-интерфейса для поиска предстоящих спортивных заездов, ознакомления с их регламентом и картографическим треком, регистрации на старты, автоматического мониторинга погодных условий на дату проведения и ведения личной спортивной статистики.

\subsubsection{Эксплуатационное назначение}
ИС «43Колеса» предназначена для эксплуатации на серверных мощностях под управлением ОС семейства Linux в контейнеризированной среде Docker с клиентским доступом через современные веб-браузеры (Chrome, Firefox, Safari, Edge, Яндекс.Браузер) на настольных компьютерах, планшетах и смартфонах.

% ==============================================================================
% 4. АНАЛИЗ ПРЕДМЕТНОЙ ОБЛАСТИ
% ==============================================================================
\section{Анализ предметной области}

\subsection{Описание предметной области}
Предметная область охватывает сферу любительского и полупрофессионального велосообщества. Ежегодно в регионе проводятся десятки разноплановых мероприятий: шоссейные критериумы, кросс-кантри марафоны, гравийные заезды и бреветы. В настоящее время информационные потоки рассредоточены по социальным сетям, что затрудняет координацию участников и сбор статистики.

\subsection{Обзор аналогов}
\begin{enumerate}
	\item \textbf{Сообщества ВКонтакте:} обеспечивают большой охват, но лишены интерактивного расписания, структурированных карточек соревнований и фильтрации по трекам;
	\item \textbf{velogorod.online (Archivelo):} содержит архивы маршрутов, однако не предоставляет инструментов регистрации на события и динамического личного кабинета;
	\item \textbf{Trello:} эффективен для организационной работы, но не приспособлен для широкого спортивного сообщества, картографии и отслеживания погоды;
	\item \textbf{MapMagic:} обеспечивает качественное построение треков, но не выполняет функций социальной платформы и афиши событий.
\end{enumerate}

% ==============================================================================
% 5. ТРЕБОВАНИЯ К РЕЗУЛЬТАТАМ РАЗРАБОТКИ
% ==============================================================================
\section{Требования к результатам разработки}

\subsection{Правила игры (бизнес-правила функционирования)}
\begin{itemize}
	\item Регистрация на событие доступна только авторизованным пользователям до момента наступления дедлайна или исчерпания лимита свободных слотов;
	\item Карточка мероприятия создается организатором со статусом «Черновик» и становится публичной на Timeline-доске только после утверждения;
	\item Изменение трека заезда организатором инициирует автоматическое оповещение всех зарегистрированных участников через Telegram/VK-бота;
	\item Завершенные события автоматически перемещаются в архив с фиксацией итогового протокола и начислением километров в профили финишировавших спортсменов.
\end{itemize}

\subsection{Требования к функциям}
Система должна обеспечивать реализацию следующих базовых подсистем:
\begin{itemize}
	\item \textbf{Подсистема авторизации и управления профилями:} регистрация по адресу электронной почты/паролю, аутентификация через JWT-токены, редактирование профиля;
	\item \textbf{Подсистема Timeline-доски мероприятий:} хронологическое отображение заездов, фильтрация по дисциплине (шоссе, МТБ, гравий), статусу и дате проведения;
	\item \textbf{Подсистема картографии:} рендеринг треков маршрута из GPX-файлов на базе OpenStreetMap/Leaflet с профилем высот;
	\item \textbf{Подсистема метеорологического мониторинга:} запрос прогноза погоды через внешний API (Open-Meteo) на координаты стартового городка на дату заезда;
	\item \textbf{Подсистема личного кабинета спортсмена:} отображение истории участий, суммарного наката, избранных заездов;
	\item \textbf{Подсистема организатора («Коллега»):} интерфейс создания события, прикрепления GPX-файла, выгрузки списка зарегистрированных спортсменов;
	\item \textbf{Подсистема уведомлений:} взаимодействие с ботом ВКонтакте для отправки сервисных сообщений и анонсов.
\end{itemize}

\subsection{Требования к показателям назначения}
\begin{itemize}
	\item Время отклика бэкенд-сервера на стандартные API-запросы: не более 1.5 секунд;
	\item Время первоначальной загрузки страниц интерфейса: не более 3.0 секунд;
	\item Пропускная способность: устойчивая работа при одновременной активности не менее 200 пользователей;
	\item Вместимость базы данных: хранение не менее 10~000 мероприятий и 50~000 профилей.
\end{itemize}

\subsection{Требования к пользовательскому интерфейсу}
\begin{itemize}
	\item Адаптивность: корректное отображение на экранах шириной от 360 px до 1920 px;
	\item Реализация темной темы по умолчанию со строгой спортивной контрастностью;
	\item Состояния интерактивных компонентов: наличие выраженных состояний покоя, наведения (hover), нажатия (active), загрузки (loading) и недоступности (disabled);
	\item Минималистичность: отсутствие перегруженных декоративных элементов, приоритизация читаемости картографических и табличных данных.
\end{itemize}

\subsection{Требования к видам обеспечения}

\subsubsection{Требования к математическому обеспечению}
Математическое обеспечение должно включать:
\begin{itemize}
	\item алгоритм вычисления длины трека по геодезическим координатам на основе ортодромической формулы гаверсинусов (Haversine formula);
	\item алгоритмы интерполяции и фильтрации шума высотного профиля маршрута;
	\item формулы расчета накопительной спортивной статистики пользователя (сумма дистанций, средняя протяженность заезда, частота стартов).
\end{itemize}

\subsubsection{Требования к информационному обеспечению}
Информационное обеспечение организуется на базе реляционной СУБД PostgreSQL.
В состав схемы данных на этапе проектирования входят ключевые сущности:
\begin{itemize}
	\item \textbf{Users} (идентификаторы, хеши паролей, email, ФИО, город, роль);
	\item \textbf{Events} (наименование, описание, дисциплина, дата/время старта, лимит мест, статус);
	\item \textbf{Routes} (геометрические полилинии, дистанция, набор высоты, ссылка на файл GPX);
	\item \textbf{EventRegistrations} (связь пользователя и события, статус заявки, дата регистрации);
	\item \textbf{WeatherCache} (координаты, дата прогноза, температура, осадки, скорость ветра).
\end{itemize}

\paragraph{5.5.2.1 Требования к форматам хранения данных}
\begin{itemize}
	\item Межсервисный обмен данными: формат JSON (JavaScript Object Notation);
	\item Геоданные и треки: форматы GPX и GeoJSON;
	\item Временные метки: международный стандарт ISO 8601 (UTC).
\end{itemize}

\paragraph{5.5.2.2 Требования к лингвистическому обеспечению}
\begin{itemize}
	\item Язык пользовательского интерфейса, системных сообщений и ТЗ: русский;
	\item Язык исходного кода, названий сущностей БД, спецификаций OpenAPI и коммитов: английский;
	\item Кодировка текстовых данных: UTF-8.
\end{itemize}

\subsubsection{Требования к метрологическому обеспечению}
\begin{itemize}
	\item Измерение дистанции маршрута: километры с округлением до сотых (0.01 км);
	\item Измерение набора высоты: метры с точностью до целых значений (1 м);
	\item Географические координаты: стандарт WGS-84 с точностью не менее 5 знаков после запятой;
	\item Фиксация времени: синхронизация системного времени сервера по протоколу NTP.
\end{itemize}

\subsubsection{Требования к техническому обеспечению}
\begin{itemize}
	\item Сервер: двухъядерный процессор архитектуры x86-64 (от 2.0 ГГц), оперативная память от 4 Гб, твердотельный накопитель SSD от 40 Гб;
	\item Клиентские устройства: персональные компьютеры и смартфоны с поддержкой стандартов HTML5, CSS3, JavaScript (ES6+).
\end{itemize}

\subsection{Требования к надежности}
\begin{itemize}
	\item Коэффициент доступности системы (uptime): не ниже 99.0\% в круглосуточном режиме;
	\item Автоматическое резервное копирование базы данных PostgreSQL: ежедневно с глубиной хранения копий 14 календарных дней;
	\item Автоматический перезапуск упавших контейнеров политиками Docker Compose (\texttt{restart: unless-stopped}).
\end{itemize}

\subsection{Требования к безопасности}
\begin{itemize}
	\item Передача данных исключительно по защищенному протоколу HTTPS с валидным SSL/TLS-сертификатом;
	\item Хеширование паролей пользователей надежными алгоритмами (bcrypt или argon2);
	\item Защита от типовых уязвимостей веб-приложений (OWASP Top 10): экранирование SQL-инъекций средствами ORM, санитизация входных данных от XSS, проверка CSRF-токенов.
\end{itemize}

\subsection{Требования к патентной чистоте}
Программные средства, библиотеки и фреймворки, используемые в проекте (Solid.JS, Python, PostgreSQL, Leaflet, Docker), должны распространяться под свободными лицензиями MIT, Apache 2.0 или BSD. Система не должна нарушать патентных прав третьих лиц.

\subsection{Требования к перспективам развития}
Архитектура системы должна предусматривать возможность последующего масштабирования:
\begin{itemize}
	\item внедрение мобильного приложения (Progressive Web App — PWA);
	\item добавление функционала трекинга участников на дистанции в реальном времени;
	\item интеграция со спортивными часами и сервисами фитнес-трекинга (Garmin Connect, Strava).
\end{itemize}

% ==============================================================================
% 6. СОСТАВ И СОДЕРЖАНИЕ РАБОТ
% ==============================================================================
\section{Состав и содержание работ}

Разработка веб-платформы «43Колеса» разделяется на следующие основные этапы:
\begin{enumerate}
	\item \textbf{Аналитический этап:} исследование предметной области, формулирование требований, составление и утверждение настоящего ТЗ.
	\item \textbf{Этап проектирования:} архитектурное моделирование системы в нотациях UML/IDEF, проектирование реляционной схемы базы данных PostgreSQL, создание прототипов пользовательского интерфейса в Figma.
	\item \textbf{Этап программной реализации:}
	      \begin{itemize}
		      \item конфигурирование окружения разработки, настройка Docker Compose;
		      \item разработка серверной части на Python (FastAPI), организация REST API;
		      \item реализация клиентского приложения на базе Solid.JS и Tailwind CSS;
		      \item интеграция интерактивной карты Leaflet и погодного API.
	      \end{itemize}
	\item \textbf{Этап тестирования:} модульное и интеграционное тестирование, проверка корректности работы алгоритмов и интерфейса.
	\item \textbf{Этап развертывания и сдачи:} публикация веб-сервиса на тестовом сервере, оформление программной документации и сдача работы.
\end{enumerate}

% ==============================================================================
% 7. ПОРЯДОК РАЗРАБОТКИ
% ==============================================================================
\section{Порядок разработки}

\begin{itemize}
	\item Разработка ведется в соответствии с гибкой методологией (Agile) и фиксацией задач в трекере;
	\item Контроль версий осуществляется в системе Git с использованием ветвления (ветка \texttt{main} — стабильная версия, \texttt{develop} — разработка, отдельные ветки \texttt{feature/*} для функционала);
	\item Каждый этап завершается код-ревью и согласованием промежуточных результатов с научным руководителем.
\end{itemize}

% ==============================================================================
% 8. ТРЕБОВАНИЯ К ДОКУМЕНТИРОВАНИЮ
% ==============================================================================
\section{Требования к документированию}

По завершении разработки исполнитель формирует комплект программных документов:
\begin{itemize}
	\item Пояснительная записка к курсовому/лабораторному проекту;
	\item Руководство оператора (пользователя) с пошаговыми инструкциями работы в веб-интерфейсе;
	\item Руководство системного администратора по развертыванию Docker-контейнеров и конфигурированию сервера;
	\item Автоматически генерируемая спецификация REST API (Swagger/OpenAPI).
\end{itemize}

% ==============================================================================
% 9. ТРЕБОВАНИЯ К ПРИЕМНО-СДАТОЧНЫМ ПРОЦЕДУРАМ
% ==============================================================================
\section{Требования к приемно-сдаточным процедурам}

Приемка разработанной веб-платформы осуществляется в соответствии с утвержденной Программой и методикой испытаний (ПМИ). Приемо-сдаточные испытания включают проверку:
\begin{itemize}
	\item успешности прохождения сценариев регистрации, авторизации и восстановления пароля;
	\item корректности вывода и фильтрации событий на Timeline-доске мероприятий;
	\item корректности отображения GPX-трека на интерактивной карте и расчета параметров дистанции;
	\item соответствия интерфейса заявленным требованиям адаптивности на десктопных и мобильных устройствах;
	\item наличия полного комплекта технической и пользовательской документации.
\end{itemize}

\end{document}